\chapter{The Problem with Time}

If the four assumptions hold, then subjective experience can emerge from static noisy information. For some reason,
this is hard to accept for many. If the computer does not run, then nothing happens. Defnitely, no Alice emerges.

\section{Alice as a Video}

Consider an ordinary particle simulation of Alice's universe. We render it frame by frame into an MPEG video. This is the straightforward, non-optimized implementation.

The ``pixels'' in this video are not flat 2D rasters, but full 3D particles. It is a complete, particle-level recording of the entire simulated universe, including Alice's body, her brain states, and every interaction. This video is an axiomatic virtual copy of real-world Alice --- functionally identical at every relevant level. Every moment of her existence, every experience, every feeling of pain or joy, is encoded in this static data.

Now we begin optimizing the code. First, we replace every \texttt{log10()} call with a precomputed lookup table. We render the simulation again. The resulting video is \emph{identical} to the original.

Next, we replace all \texttt{sqrt()} computations with precomputed values. Again, the videos are indistinguishable.

We continue this process, gradually replacing more and more active computation with static lookup tables. At every step, the output video remains exactly the same. The only difference is the time required to generate it: the more we optimize, the faster the rendering completes.

In the end, we optimize the simulation completely. All computation has been replaced by precomputed data. What remains is pure static information --- the video file itself.

All these videos are indistinguishable from one another. Only the code/data ratio of the process that created them differs.

This raises a direct question: Is some minimum level of active computation really required to maintain Alice's subjective experience? If so, what is the critical code/data ratio below which consciousness disappears?

If no such threshold exists, then subjective experience can persist in purely static information.

If subjective experience gradually faded with the optimizations, then there would be a direct contradiction with the Axiom-4, that pain affects Alices behavior. All the videos are equal and
demonstrate identical behavior.


\section{The Self-Mutating Computer}

Consider the following thought experiment.

Imagine a computer running a high-fidelity DNA simulation of Alice. Now suppose we add a special ``self-mutating'' function to its code.
Instead of terminating after one run, this function modifies the computer's own memory and instructions, then executes the new configuration. This process repeats indefinitely.

The computer systematically explores \emph{every possible way} to interpret its finite number of bits $n$. It tries every possible division between code and data, every possible ordering of execution blocks, every possible endianness, and every possible interpretation of the bit strings. 

Since there are $n$ bits, there exist $2^n$ possible configurations. When we also consider all possible orderings (permutations), repetitions, and sub-paths, the machine will eventually execute \emph{every possible program} that can be encoded within those $n$ bits.

Among this enormous space of executed programs, most will produce nothing but noise. However, a tiny subset will produce highly compressible, coherent structures --- including stable geometries, consistent physical laws, and conscious observers such as Alice.

In this picture, there is no need for an external ``reader'' or simulator. The universe is self-contained: all possible interpretations and all possible observers exist as different ways the same information can be organized and traversed.

One might object that this process still depends on an external clock to drive the mutations and executions. However, this objection is easily overcome. The ``clock'' itself can be internalized --- made part of the information and the self-mutating process. Once the clock is moved inside the system, the entire construction becomes completely self-contained. There is no longer any need for anything external. The information generates its own dynamics, its own time, and its own observers.

This self-mutating computer provides an intuitive bridge toward understanding how entire universes --- complete with conscious beings experiencing geometric reality, pain, joy, and the flow of time --- can arise purely from information without any outside intervention.

Moving the external clock inside the computer, turns  the computer from runnign to static ``block universe''.

In set-theoretic language, we have to do is to embed the computer and the software it is running to a super-set of both. Now we have
a staitc set, in which ``running'' is just one structure with its interpretation.


